\chapter{Data and Empirical Strategy}\label{ch:data}

\section{Data and Summary Statistics}\label{sec:data}

The national data are a crop-year panel of 43 crop or crop-form series from the Crop Production Survey \citep{statkorea_nd}, analyzed over 1991-2024. Cultivated area is in hectares, yield in kilograms per 10 ares, and production in metric tonnes. One hectare contains ten units of 10 ares, and one tonne contains 1,000 kilograms. Thus production equals area \texttimes{} yield \texttimes{} 10/1,000, or area \texttimes{} yield / 100, up to reporting and rounding differences. Equivalently, 100 kg/10a equals 1 tonne/ha. With identical observations and weights, the log production contrast is approximately the sum of the log area and yield contrasts. This is an accounting decomposition; average yield can reflect both productivity and the composition of cultivated land.

Exposure is defined by the availability of insurance, not by enrollment, and national totals include insured and uninsured farms. Two changes in survey methods affect comparability: sweet potato moved to administrative reporting in 2005, and spring napa cabbage and spring radish moved to a sample survey in 2012 \citep{statkorea2014}.

Table~\ref{tab:summary} reports the mean, standard deviation, and range of the unique crop-year observations entering the main comparison. Treated crops are on average about three times larger in cultivated area. These statistics describe levels and dispersion. Differences in observed pre-pilot changes are examined in Section~\ref{sec:comparability}; similar means or a log transformation would not establish parallel untreated trends.

\begin{table}[htbp]
\centering
\caption{Summary Statistics}\label{tab:summary}
\begin{tabular*}{\textwidth}{@{\extracolsep{\fill}}lrrrrcc}
\toprule
 & Mean & S.D. & Min. & Max. & Obs. & Crops \\
\midrule
\multicolumn{7}{@{}l}{\textit{Panel A. Area (ha)}} \\
Treated, reference year     & 31,958 & 24,132 & 13,033 & 76,267 & 7   & 7  \\
Treated, target years       & 26,991 & 15,482 & 13,017 & 80,842 & 70  & 7  \\
Comparison, reference years & 9,775  & 9,659  & 1,062  & 34,875 & 66  & 22 \\
Comparison, target years    & 9,083  & 10,014 & 619    & 45,474 & 218 & 22 \\[4pt]
\multicolumn{7}{@{}l}{\textit{Panel B. Yield (kg/10a)}} \\
Treated, reference year     & 1,934 & 2,311 & 150 & 6,725  & 7   & 7  \\
Treated, target years       & 1,856 & 2,161 & 147 & 8,541  & 70  & 7  \\
Comparison, reference years & 2,488 & 2,696 & 37  & 10,946 & 66  & 22 \\
Comparison, target years    & 2,335 & 2,757 & 30  & 11,286 & 218 & 22 \\[4pt]
\multicolumn{7}{@{}l}{\textit{Panel C. Production (t)}} \\
Treated, reference year     & 350,922 & 323,462 & 92,830 & 1,035,076 & 7   & 7  \\
Treated, target years       & 381,113 & 424,756 & 55,714 & 1,594,450 & 70  & 7  \\
Comparison, reference years & 208,852 & 330,283 & 1,589  & 1,582,966 & 66  & 22 \\
Comparison, target years    & 163,544 & 295,436 & 1,598  & 1,698,463 & 218 & 22 \\
\bottomrule
\end{tabular*}
\notes{Statistics are computed over unique crop-year observations that enter the main comparison. S.D. is the standard deviation of observations. Reference years are crop-specific; target years are the ten years beginning with the coded expansion year.}
\end{table}

\section{Sample and Exposure Timing}\label{sec:sample}

The treated sample consists of seven annual crops (soybean, onion, sweet potato, corn, garlic, spring potato, and red pepper), each with a continuous series from 1991, a dated pilot, and ten observed years after expansion. Fruit crops are excluded because orchard area and perennial production respond on different time scales, and autumn potato is excluded because its milestones are not reported separately from other potato forms. Rice is analyzed separately through its municipal rollout (Section~\ref{sec:rice_design}).

Table~\ref{tab:timing} reports the timing used to organize the comparisons. The pilot year marks regional availability, and the coded expansion year begins the ten target years. Neither implies universal enrollment or geographical coverage. Historical announcement dates are converted into harvest years using recent sales calendars, creating timing uncertainty. Spring potato remained restricted to selected regions in 2023-2024 \citep{mafra2024,mafra2025}, so its coded expansion year should not be interpreted as nationwide availability. Alternative milestone definitions are examined in Table~\ref{tab:sensitivity}.

\begin{table}[htbp]
\centering
\caption{Insurance Timing for the Seven Annual Crops}\label{tab:timing}
\begin{tabular*}{\textwidth}{@{\extracolsep{\fill}}lccccc}
\toprule
Crop & \shortstack{Pilot\\year} & \shortstack{Expansion\\year} & \shortstack{Product-table\\year} & \shortstack{Reference\\year} & \shortstack{Transition\\years} \\
\midrule
Soybean       & 2008 & 2012 & 2011 & 2007 & 4 \\
Onion         & 2009 & 2012 & 2011 & 2008 & 3 \\
Sweet potato  & 2009 & 2013 & 2012 & 2008 & 4 \\
Corn          & 2009 & 2013 & 2012 & 2008 & 4 \\
Garlic        & 2010 & 2013 & 2012 & 2009 & 3 \\
Spring potato & 2008 & 2013 & 2011 & 2007 & 5 \\
Red pepper    & 2008 & 2015 & 2014 & 2007 & 7 \\
\bottomrule
\end{tabular*}
\notes{Pilot and expansion years follow the milestone sequence described in Section~\ref{sec:expansion}. Reference year is the pilot year minus one. Transition years run from the pilot year to the year before expansion. Product-table years are an alternative milestone used in Table~\ref{tab:sensitivity}. Data from \citet[pp.~37, 44, 46]{mafra2025}.}
\end{table}

The candidate comparison pool contains 24 annual crop series, of which 22 contribute to the main contrast. A crop is eligible as a comparison only if both its reference and target observations precede its own pilot, or if it has no pilot by 2024. Because all seven treated crops entered pilots by 2010, before the earliest target year of 2012, they never serve as comparisons for one another in the main post-period contrasts.

\section{National Crop Comparisons}\label{sec:national_design}

The national analysis constructs crop-specific difference-in-differences contrasts and averages them across seven crops and ten target years. Let $p_{i}$ be the pilot year, $g_{i}$ the expansion year, and $b_{i} = p_{i} - 1$ the reference year of treated crop $i$. At event year $e$, the target year is $g_{i} + e$, and the comparison set $C_{i,e}$ contains the $n_{i,e}$ crops observed and not yet exposed in both $b_{i}$ and $g_{i} + e$. For a log outcome $y$, the crop-event contrast is
\begin{equation}\label{eq:contrast}
\widehat{\delta}_{i,e} = \left( y_{i,g_{i} + e} - y_{i,b_{i}} \right) - \frac{1}{n_{i,e}}\sum_{j \in C_{i,e}}\left( y_{j,g_{i} + e} - y_{j,b_{i}} \right)
\end{equation}
and the contrasts are averaged with equal weights across crops and event years:
\begin{equation}\label{eq:average}
\widehat{\theta}_{e} = \frac{1}{7}\sum_{i = 1}^{7}\widehat{\delta}_{i,e},\qquad \widehat{\theta} = \frac{1}{10}\sum_{e = 0}^{9}\widehat{\theta}_{e}.
\end{equation}

The first term in equation~\eqref{eq:contrast} is the log growth of the treated crop from its reference year, and the second is the average log growth of its comparison crops over the same interval. Each contrast is therefore a difference in log changes from the reference year: the growth of the treated crop in excess of the growth of crops not yet insured. The statistic $\widehat{\theta}$ averages these net growth rates over seven crops and ten target years, and multiplied by 100 it approximates a difference in percentage growth. Because the reference year precedes the pilot, $\widehat{\theta}$ covers the pilot transition as well as the period after expansion.

As an alternative aggregation, crops are weighted by their mean cultivated area in 2003-2007, before the earliest pilot:
\begin{equation}\label{eq:areaweights}
\widehat{\theta}^{A} = \frac{1}{10}\sum_{e = 0}^{9}\sum_{i = 1}^{7}w_{i}^{A}\,\widehat{\delta}_{i,e},\qquad w_{i}^{A} = \frac{\bar{A}_{i}^{pre}}{\sum_{k = 1}^{7}\bar{A}_{k}^{pre}}
\end{equation}

Event-time contrasts use the pilot year to locate observations on the horizontal axis while retaining the last pre-pilot year as the reference observation. For each event year, equation~\eqref{eq:contrast} compares the treated crop's change with that of eligible comparison crops, and the seven contrasts are averaged equally. The last pre-pilot contrast is zero by construction. Earlier contrasts describe observed pre-pilot differences; their joint test assesses whether nine contrasts are zero. Neither a small contrast nor a non-rejection establishes parallel untreated trends. Post-pilot contrasts include the transition before coded expansion.

\section{Municipal Rice Comparison}\label{sec:rice_design}

Rice insurance was first offered in 20 municipalities in 2009, with ten further municipalities added in 2011 \citep{choi2010,mifaff2011}. Both cohorts were designated as major producing areas. Later entrants provide a comparison for early entrants through 2010, before their own eligibility. A shared designation holds some features of program selection constant but does not make rollout timing random. Section~\ref{sec:rice_results} compares observed pre-pilot characteristics and paths.

The outcome is log paddy-rice planted area for 2000-2010 from the Crop Production Survey. Restricting the sample to mainland municipalities with complete, boundary-consistent series leaves 17 early and nine later entrants, or 286 municipality-years. The estimating equation is
\begin{equation}\label{eq:rice}
\ln A_{mt} = \alpha_{m} + \lambda_{t} + \beta D_{mt} + \varepsilon_{mt}
\end{equation}
where $A_{mt}$ is planted area in municipality $m$ and year $t$, $\alpha_{m}$ and $\lambda_{t}$ are municipality and year fixed effects, and $D_{mt}$ equals one for early entrants in 2009-2010. Municipality fixed effects absorb time-invariant differences such as soil quality, irrigation infrastructure, and proximity to markets; year fixed effects absorb national shocks such as rice prices and weather common to all municipalities.

\section{Identification, Inference, and Scope}\label{sec:identification}

A causal interpretation requires that treated and comparison units would have followed parallel changes without insurance, with correctly dated exposure and no anticipation or spillovers. For the national comparison, crop substitution or common-market responses could affect comparison crops even before their own eligibility. Pre-pilot paths and group characteristics assess observed comparability but cannot establish the untreated post-pilot path. The rice comparison holds crop identity fixed; regional weather, investment, and other time-varying local conditions may nevertheless differ between cohorts.

National uncertainty is summarized with crop-level scores and independent Webb six-point multipliers (Appendix~\ref{app:inference}). The procedure is neither null-imposed nor studentized, and its small-sample coverage is not established. Webb weights alone do not resolve few-treated-cluster problems \citep{mackinnonwebb2018}. Municipal inference uses a null-imposed, CR1-studentized wild cluster bootstrap with Webb weights and 9,999 draws, with intervals obtained by test inversion \citep{cameron2008,webb2023}. Finite-sample concerns remain with 26 municipality clusters, including nine comparisons. Tables report standard errors in parentheses, 95 percent intervals in brackets, and * for p < 0.10, ** for p < 0.05, and *** for p < 0.01. Statistical significance otherwise refers to the 5 percent level.

Uncertainty and identification are separate considerations. A wide interval limits precision even under valid identifying assumptions; a narrow interval cannot resolve a biased comparison.\footnote{A normal-approximation benchmark, treating the reported standard errors as fixed (0.159 for the national production contrast and 0.022 for the rice estimate), gives minimum detectable effects of approximately 0.45 and 0.06 log points for 80 percent power at a two-sided 5 percent level. These calculations do not establish the power of the bootstrap procedures used here.} The national intervals are exploratory because the score procedure's nominal coverage is uncertain. Alternative reference periods, comparison pools, and weights were examined after the main results were known and are reported as sensitivity analyses.
